% ###################################################################################################
% Chapter 6 - Clothing and Equipment
% ###################################################################################################
\chapter{Clothing and Equipment}

In Eidos, what characters wear is not just a matter of protection—it reflects their background, status, and the role they play in the world. A noble dressed in ceremonial robes conveys a different message than a seasoned warrior in battle-worn mail. Clothing and armor help shape the narrative, affect how NPCs react, and influence a character’s abilities within the game. Players are encouraged to select outfits that align with their character’s role, culture, and personal history.

Each character will have at least two sets of clothing: one primary and one secondary. These sets represent their lifestyle and choices. For example, a merchant might carry both casual clothing for travel and finer garments for negotiations, while a warrior may have ceremonial armor for formal events and heavier armor for battle. However, the use of mismatched or inappropriate gear—for example, wearing armor in court or dressing poorly for a formal audience—can impact interactions and story outcomes.

This system emphasizes that clothing and armor aren’t just mechanical choices but integral to role-playing. What characters wear shapes their place in the world, reflects their status, and influences how others perceive them. The chosen attire should make sense given their station, environment, and circumstances.

\section{Armor Levels and Narrative Effects}

Each set of clothing or armor falls into one of four levels: No Armor, Lightly Armored, Armored, or Heavily Armored. Each level brings specific benefits and limitations, encouraging players to think strategically about what to wear in different scenarios.

\subsection*{No Armor}

\textbf{Description:} Simple clothes with no protective qualities—robes, tunics, or travel gear.

\begin{itemize}
    \item \textbf{Pros:}
    \begin{itemize}
        \item Maximum freedom of movement and agility.
        \item Ideal for stealth and blending into crowds.
        \item No exhaustion from prolonged wear.
    \end{itemize}
    \item \textbf{Cons:}
    \begin{itemize}
        \item Vulnerable to all forms of physical attacks.
        \item Requires reliance on clever tactics or diplomacy in dangerous situations.
    \end{itemize}
\end{itemize}

\subsection*{Lightly Armored}

\textbf{Description:} Leather armor, padded jackets, or chainmail worn under clothing—suitable for agile fighters or scouts.

\begin{itemize}
    \item \textbf{Pros:}
    \begin{itemize}
        \item Offers moderate protection without sacrificing much agility.
        \item Useful for roles that demand both defense and speed.
    \end{itemize}
    \item \textbf{Cons:}
    \begin{itemize}
        \item Can’t stop heavier strikes or blunt force.
        \item Mild discomfort when worn for long periods, especially in hot climates.
    \end{itemize}
\end{itemize}

\subsection*{Armored}

\textbf{Description:} Chainmail or partial plate—common for soldiers or knights who need balanced defense and movement.

\begin{itemize}
    \item \textbf{Pros:}
    \begin{itemize}
        \item Protects against slashing and piercing attacks.
        \item Intimidating to opponents and suitable for frontline combat.
    \end{itemize}
    \item \textbf{Cons:}
    \begin{itemize}
        \item Noisy and limits stealth attempts.
        \item Slows down the wearer during long chases or escapes.
    \end{itemize}
\end{itemize}

\subsection*{Heavily Armored}

\textbf{Description:} Full plate armor or reinforced chainmail, designed for the battlefield.

\begin{itemize}
    \item \textbf{Pros:}
    \begin{itemize}
        \item Provides the highest level of protection, resisting most physical damage.
        \item Creates an overwhelming presence, discouraging enemies.
    \end{itemize}
    \item \textbf{Cons:}
    \begin{itemize}
        \item Severely limits movement—running, climbing, or swimming is nearly impossible.
        \item Causes exhaustion if worn for too long or in the wrong environment.
    \end{itemize}
\end{itemize}

\section{Changing and Using Armor Narratively}

The choice of clothing and armor is not permanent—characters can change between their sets as needed, but doing so takes time and planning. For instance:

\begin{itemize}
    \item Swapping from travel clothes to ceremonial robes before entering a royal audience can influence the outcome of a negotiation.
    \item Switching from light to heavy armor before a siege offers protection but limits agility.
    \item Choosing not to wear armor at all during a diplomatic event signals trust but also leaves the character vulnerable if things go wrong.
\end{itemize}

Armor not only defines how characters look but also shapes the gameplay. A heavily armored warrior might struggle to infiltrate a guarded mansion, while a rogue in light clothing could breeze through. Each level comes with strengths and weaknesses that push players to think creatively, adapting their equipment to suit the challenges of the story.

\section{Active Armor Mechanics: Severity Reduction}

While armor offers clear narrative advantages, it also possesses a vital mechanical function in combat. In Eidos, armor does not make you harder to hit (which is governed by Agility and shields). Instead, armor directly absorbs the kinetic energy, claws, or blades of a strike, reducing the severity of incoming wounds.

When a character is hit and the attacker rolls on the Critical Injury tables (Chapter 7), the defender subtracts their armor's \textbf{Severity Reduction} value directly from the attacker's $1\text{d}20$ injury roll:

\begin{center}
\begin{tabular}{lcp{8cm}}
\hline
\textbf{Armor Level} & \textbf{Severity Reduction} & \textbf{Typical Examples} \\ \hline
No Armor & -0 & Silk robes, linen tunics, heavy cloaks. \\
Lightly Armored & -2 & Padded Gambeson, hardened leather, light brigandine. \\
Armored & -4 & Iron chainmail, partial steel plate, ringmail. \\
Heavily Armored & -6 & Full plate harness, reinforced heavy scale. \\ \hline
\end{tabular}
\end{center}

\textbf{Superficial Scratches:} If the Severity Reduction modifies the $1\text{d}20$ critical injury roll to \textbf{0 or lower}, the strike is fully absorbed or deflected by the armor's integrity. The defender suffers only a superficial scratch, resulting in a minor bruise requiring 1 hour of recovery, with no lasting penalties or permanent scars.

\section{Weapons and Combat Gear}

Weapons are categorized by weight and size, determining which core Attribute is used to roll attack checks (Chapter 7) and what damage type is inflicted. The three damage types—\textbf{Cutting (C)}, \textbf{Stabbing (S)}, and \textbf{Impact (I)}—correspond directly to the critical injury tables in Chapter 7.

\subsection{Weapons Table}

\begin{center}
\begin{tabular}{llcccc}
\hline
\textbf{Weapon Name} & \textbf{Category} & \textbf{Damage Type} & \textbf{Hands} & \textbf{Cost} & \textbf{Weight} \\ \hline
Dagger & Light (AGI) & Stabbing (S) & 1 & 2 gp & 1 lb \\
Throwing Knife & Light (AGI) & Stabbing (S) & 1 & 1 gp & 0.5 lb \\
Shortsword & Light (AGI) & Cutting (C) & 1 & 8 gp & 2 lbs \\
Rapier & Light (AGI) & Stabbing (S) & 1 & 15 gp & 2 lbs \\
Arming Sword & Medium (AGI/STR) & Cutting (C) / Stabbing (S) & 1 & 12 gp & 3 lbs \\
Hand Axe & Medium (AGI/STR) & Cutting (C) & 1 & 4 gp & 3 lbs \\
Mace & Medium (AGI/STR) & Impact (I) & 1 & 6 gp & 4 lbs \\
Spear & Medium (AGI/STR) & Stabbing (S) & 1 or 2 & 3 gp & 4 lbs \\
Greatsword & Heavy (STR) & Cutting (C) & 2 & 30 gp & 7 lbs \\
War Hammer & Heavy (STR) & Impact (I) & 2 & 20 gp & 6 lbs \\
Battleaxe & Heavy (STR) & Cutting (C) & 2 & 15 gp & 6 lbs \\
Halberd & Heavy (STR) & Cutting (C) / Stabbing (S) & 2 & 25 gp & 8 lbs \\
Shortbow & Ranged (AGI) & Stabbing (S) & 2 & 10 gp & 2 lbs \\
Longbow & Ranged (AGI) & Stabbing (S) & 2 & 25 gp & 3 lbs \\ \hline
\end{tabular}
\end{center}

\subsection{Shields}

Shields are vital active and passive defense tools, allowing characters to block strikes and increase their survivability.

\begin{center}
\begin{tabular}{lcccr}
\hline
\textbf{Shield Type} & \textbf{Passive Defense Bonus} & \textbf{Active Block Modifier} & \textbf{Cost} & \textbf{Weight} \\ \hline
Buckler & +1 & +2 & 3 gp & 2 lbs \\
Round Shield & +2 & +1 & 7 gp & 5 lbs \\
Tower Shield & +4 & -1 & 15 gp & 12 lbs \\ \hline
\end{tabular}
\end{center}

\begin{itemize}
    \item \textbf{Passive Defense Bonus:} Added directly to the static Defense rating ($10 + \text{AGI} + \text{Shield Bonus}$).
    \item \textbf{Active Block Modifier:} Added to Active Defense rolls when declaring a shield block reaction. A larger shield is heavier and slower to swing into position for a parry, but provides wider passive cover.
\end{itemize}

\section{Economy and Adventuring Gear}

The economy of Eidos is based on pre-modern commerce, relying on coinage minted from copper, silver, and gold.

\begin{center}
    \textbf{1 Gold Piece (gp) = 10 Silver Pieces (sp) = 100 Copper Pieces (cp)}
\end{center}

\subsection{Adventuring Gear Table}

\begin{center}
\begin{tabular}{lccr}
\hline
\textbf{Item} & \textbf{Description} & \textbf{Cost} & \textbf{Weight} \\ \hline
Backpack & Holds up to 40 lbs of gear & 2 gp & 2 lbs \\
Rations (1 day) & Dry meats, hard tack, nuts & 5 sp & 1 lb \\
Waterskin & Holds 2 pints of water & 1 gp & 4 lbs (full) \\
Torch & Burns for 1 hour, sheds light & 1 sp & 1 lb \\
Flint and Steel & Used to light fires & 1 gp & - \\
Hemp Rope (50 ft) & Sturdy rope for climbing & 1 gp & 10 lbs \\
Bedroll & A padded canvas roll for sleeping & 1 gp & 5 lbs \\
First Aid Kit & Bandages, herbs, splints (5 uses) & 5 gp & 3 lbs \\
Lantern & Sheds light, requires oil & 8 gp & 2 lbs \\
Oil Flask (1 pint) & Fuel for lantern or starting fires & 1 sp & 1 lb \\ \hline
\end{tabular}
\end{center}

